\documentclass[11pt,a4paper]{article}
\usepackage{fontspec}
\setmainfont{Latin Modern Roman}
\usepackage{amsmath,amssymb,amsthm}
\usepackage[margin=1in]{geometry}
\usepackage{hyperref}
\usepackage{booktabs}
\hypersetup{colorlinks=true,linkcolor=blue,citecolor=blue,urlcolor=blue}

\newtheorem{theorem}{Theorem}
\newtheorem{proposition}[theorem]{Proposition}
\newtheorem{corollary}[theorem]{Corollary}
\newtheorem*{remark}{Remark}

\title{\textbf{Null congruence self-consistency as a characterisation\\
of the Kerr vacuum in the principal null direction}}
\author{Y.\,Y.\,N.\ Li\footnote{Independent Researcher,
China. E-mail: \texttt{papaso@icloud.com}.}}
\date{}

\begin{document}
\maketitle

\begin{abstract}
We study self-consistency functionals $K_i:=r^2\,\Box\ln\sigma_i$,
where $\sigma_1:=r\sqrt{f}$, $\sigma_2:=r$ are scalar functions
of the static metric
$ds^2=-f\,dt^2+f^{-1}dr^2+r^2\,d\Omega^2$
and $\Box:=r^{-2}\partial_r(r^2 f\,\partial_r)$.
We prove two results.

First, for $f(r)=1-2m(r)/r$ with arbitrary $C^2$ mass function
$m(r)$, three exact identities hold:
$K_1=1-\tfrac{1}{2}r^2R$,\;$K_2=1-8\pi r^2\rho$,\;
$K_1-K_2=8\pi r^2 p_{\!\perp}$,
so $\{K_1,K_2\}$ encodes all independent components of the
spherically symmetric Einstein equations.
The condition $K_1=1$ is equivalent to $r^4\rho=\mathrm{const}$,
satisfied by vacuum and by the Maxwell field
$(\rho=Q^2/(8\pi r^4))$, but not by the cosmological constant.

Second, for the Kerr metric we set $\sigma:=\sqrt{\Sigma}$
(so $\sigma_1=\sigma_2=\sigma$ and $K_1=K_2=:K$;
$\Sigma:=r^2+a^2\cos^2\!\theta$) and define
$\Box_{\mathrm{null}}\phi:=\partial_r^2\phi+(2r/\Sigma)\,\partial_r\phi$
along the principal outgoing null congruence.
Then $\Sigma\,\Box_{\mathrm{null}}\ln\sigma=1$ holds identically
for all $(r,\theta,M,a)$; the proof is three lines of algebra.
Via the Raychaudhuri equation, this condition is equivalent to
$R_{\mu\nu}\ell^\mu\ell^\nu=0$, characterising the vacuum
condition along the principal null direction.
The $\theta$-dependence of $\sqrt{\Sigma}$ is essential: no
purely radial function satisfies the self-consistency condition.
\end{abstract}

\medskip
\noindent\textbf{Keywords:} self-consistency $\cdot$ null geodesic congruence
$\cdot$ Kerr metric $\cdot$ Einstein equations
$\cdot$ Raychaudhuri equation $\cdot$ Newman--Penrose formalism

%───────────────────────────────────────────────────────
\section{Setup and notation}

Consider a static, spherically symmetric metric
\begin{equation}\label{eq:metric}
  ds^2 = -f(r)\,dt^2 + f(r)^{-1}dr^2 + r^2\,d\Omega^2,
  \qquad f(r)=1-\frac{2m(r)}{r},
\end{equation}
with arbitrary $C^2$ mass function $m(r)$.
Set $\Delta(r):=r^2 f(r)=r^2-2m(r)r$, define the symplectic
eigenvalues~\cite{K1} $\sigma_1:=r\sqrt{f}=\sqrt{\Delta}$ and $\sigma_2:=r$,
and the radial d'Alembertian and self-consistency functionals
\begin{equation}
  \Box\phi:=\frac{1}{r^2}\partial_r(r^2 f\,\partial_r\phi),
  \qquad
  K_i:=r^2\,\Box\ln\sigma_i.
\end{equation}

For the Kerr metric in Boyer--Lindquist coordinates we write
$\Sigma:=r^2+a^2\cos^2\!\theta$,\;
$\Delta:=r^2-2Mr+a^2$ (so $a=0$ reduces to the Schwarzschild
case), and work with the principal outgoing null congruence
$\ell^\mu=\bigl((r^2+a^2)/\Delta,\,1,\,0,\,a/\Delta\bigr)$
(Kinnersley form, affinely parameterised so that $\ell^r=1$,
$\ell^\theta=0$).
Its expansion is given by the Newman--Penrose spin coefficient
$\rho_{\scriptscriptstyle\mathrm{NP}}=-1/(r-ia\cos\theta)$
\cite[Sec.~56]{Chandrasekhar1983}:
\begin{equation}\label{eq:theta}
  \theta_{\mathrm{exp}}
  := -2\,\mathrm{Re}(\rho_{\scriptscriptstyle\mathrm{NP}})
  = \frac{2r}{\Sigma}.
\end{equation}
We define the null d'Alembertian for $r$-dependent functions along
this congruence:
\begin{equation}\label{eq:boxnull}
  \Box_{\mathrm{null}}\,\phi
  := \partial_r^2\phi + \theta_{\mathrm{exp}}\,\partial_r\phi.
\end{equation}
(Since $D\phi=\ell^\nu\partial_\nu\phi$ is a scalar,
$D^2\phi=\ell^\mu\partial_\mu(D\phi)$ with no connection correction;
\eqref{eq:boxnull} equals $D^2\phi+\theta_{\mathrm{exp}}\,D\phi$.)

%───────────────────────────────────────────────────────
\section{Spherically symmetric case}

\begin{theorem}[Exact curvature identities]\label{thm:spherical}
For metric~\eqref{eq:metric} with $f=1-2m(r)/r$,
the following hold identically for any $C^2$ function $m(r)$:
\begin{align}
  K_1 &= 1 - \tfrac{1}{2}r^2 R, \label{eq:K1R}\\
  K_2 &= 1 - 8\pi r^2\rho,       \label{eq:K2rho}\\
  K_1-K_2 &= 8\pi r^2 p_{\!\perp},  \label{eq:K12}
\end{align}
where $R$ is the Ricci scalar, $8\pi\rho=G^t{}_t=2m'(r)/r^2$
the energy density, and $8\pi p_{\!\perp}=G_{\theta\theta}/r^2$
the tangential pressure.
\end{theorem}

\begin{proof}
Write $\Delta=r^2-2m(r)r$; then $K_1=\Delta''/2=f+2rf'+\tfrac{1}{2}r^2f''$ and
$K_2=d(rf)/dr=f+rf'=1-2m'(r)$; the identities then follow by
comparing with the standard spherically symmetric Ricci
scalar $R=-f''-4f'/r+2(1-f)/r^2$ and Einstein components.
Equation~\eqref{eq:K12} is $G_{\theta\theta}=K_1-K_2$,
verified by $G_{\theta\theta}=\tfrac{1}{2}r^2f''+rf'$.\qed
\end{proof}

\begin{corollary}[Complete Einstein encoding]\label{cor:complete}
The pair $(K_1,K_2)$ encodes all independent components of the
spherically symmetric Einstein equations:

\smallskip
\begin{center}
\renewcommand{\arraystretch}{1.3}
\begin{tabular}{@{}lll@{}}
\toprule
Einstein component & $K_i$ expression & Content \\
\midrule
$G^t{}_t = 8\pi\rho$ &
  $K_2 = 1 - 8\pi r^2\rho$ & energy density \\
$G_{\theta\theta} = 8\pi r^2 p_{\!\perp}$ &
  $K_1-K_2 = 8\pi r^2 p_{\!\perp}$ & tangential pressure \\
\bottomrule
\end{tabular}
\end{center}
\smallskip

\noindent
The metric~\eqref{eq:metric} satisfies $g_{tt}=-g_{rr}^{-1}$, which
forces $G^r{}_r = G^t{}_t = 2m'/r^2$; the radial pressure is
therefore not independent of the energy density in these coordinates.
The two independent components are $G^t{}_t$ (encoded by $K_2$) and
$G_{\theta\theta}$ (encoded by $K_1-K_2$).

Vacuum ($\rho=p_{\!\perp}=0$) gives $K_1=K_2=1$, recovering
the vacuum equivalence $K_1=K_2=1\Leftrightarrow R_{\mu\nu}=0$
of~\cite{K1}.
\end{corollary}

\noindent\textbf{Physical interpretation of $K_1$.}\;
The exact formula $K_1=1-rm''(r)-2m'(r)$ gives
\begin{equation}\label{eq:K1cond}
  K_1=1
  \;\Leftrightarrow\;
  rm''+2m'=0
  \;\Leftrightarrow\;
  \frac{d}{dr}(r^2 m')=0
  \;\Leftrightarrow\;
  r^4\rho = \text{const}.
\end{equation}
The condition $r^4\rho=\text{const}$ holds for $\rho=0$ (vacuum),
for $\rho=Q^2/(8\pi r^4)$ (Maxwell field, Reissner--Nordstr\"{o}m),
and more generally for any $1/r^4$ energy density.
It fails for $\rho=\Lambda/(8\pi)$ (cosmological constant), which
gives $K_1-1=-2\Lambda r^2$, and for any extended mass distribution
with $\rho\propto r^n$, $n\neq -4$.
For Schwarzschild--de~Sitter, $K_1-1=-2\Lambda r^2$ and
$K_2-1=-\Lambda r^2$, recovering the effective energy density
$\rho_\Lambda=\Lambda/(8\pi)$.

%───────────────────────────────────────────────────────
\section{Kerr extension via null geodesic congruence}

The spherically symmetric result used $\sigma_2=r$, which is
only the area radius for spherically symmetric metrics.
For Kerr, the natural area function at angle $\theta$ is
$\sqrt{\Sigma}$, and the expansion $\theta_{\mathrm{exp}}$
replaces $2/r$ by $2r/\Sigma$.

\begin{theorem}[$K{=}1$ for Kerr]\label{thm:kerr}
For the Kerr metric, set $\sigma:=\sqrt{\Sigma}$ and define
$\Box_{\mathrm{null}}$ as in~\eqref{eq:boxnull}--\eqref{eq:theta}.
Then
\begin{equation}\label{eq:kerr}
  \Sigma\,\Box_{\mathrm{null}}\ln\sigma = 1
\end{equation}
holds identically for all $r,\theta,M,a$.
\end{theorem}

\begin{proof}
Three lines. With $\sigma=\sqrt{\Sigma}$:
\begin{align*}
  \partial_r\ln\sigma &= \frac{r}{\Sigma},\\
  \partial_r^2\ln\sigma &= \frac{\Sigma-2r^2}{\Sigma^2}
    = \frac{a^2\cos^2\!\theta-r^2}{\Sigma^2},\\
  \Box_{\mathrm{null}}\ln\sigma
    &= \frac{a^2\cos^2\!\theta-r^2}{\Sigma^2}
      +\frac{2r}{\Sigma}\cdot\frac{r}{\Sigma}
    = \frac{a^2\cos^2\!\theta-r^2+2r^2}{\Sigma^2}
    = \frac{\Sigma}{\Sigma^2}
    = \frac{1}{\Sigma}.
\end{align*}
Hence $\Sigma\cdot\Box_{\mathrm{null}}\ln\sigma
      =\Sigma\cdot(1/\Sigma)=1$.\qed
\end{proof}

\begin{remark}[Schwarzschild limit]
Setting $a=0$ gives $\Sigma=r^2$, $\theta_{\mathrm{exp}}=2/r$,
and $\sigma=r$.  Equation~\eqref{eq:kerr} reduces to
$r^2\cdot\Box_{\mathrm{null}}\ln r=1$, which holds for any
$m(r)$ since $\partial_r^2\ln r+(2/r)\partial_r\ln r=1/r^2$,
and recovers $K_2=1$ in the spherically symmetric notation~\cite{K1}.
The replacement $r\to\sqrt{\Sigma}$ is the unique extension
consistent with Kerr's geometry: $\sqrt{\Sigma}$ is the
transverse length scale that appears in the Newman--Penrose
expansion $\theta_{\mathrm{exp}}=2r/\Sigma$ and in the
area element of the 2-spheres cross-sectioning the null
congruence.
\end{remark}

\begin{remark}[Single eigenvalue]
In the Kerr case, $\sigma_1=\sigma_2=\sqrt{\Sigma}$: the
framework has a single eigenvalue rather than two.
This reflects the fact that the principal null congruence
is shear-free ($\sigma_{\scriptscriptstyle\mathrm{NP}}=0$;
Kerr is algebraically special, Petrov type~D), so there is
no independent transverse shape eigenvalue to distinguish.
\end{remark}

\begin{remark}[Role of $\Sigma$]
The function $\Sigma=r^2+a^2\cos^2\!\theta$ is the natural
length scale of Kerr: it governs Carter separability, appears in
every Kerr Christoffel symbol, and determines the Hawking
temperature via the surface gravity.
Theorem~\ref{thm:kerr} identifies $\sqrt{\Sigma}$ as the
eigenvalue whose self-consistency — $\sigma^2\Box_\mathrm{null}
\ln\sigma=1$ — is exact throughout the spacetime, not merely
near the horizon.
\end{remark}

Table~\ref{tab:comparison} summarises the two frameworks.

\begin{table}[h]
\centering
\renewcommand{\arraystretch}{1.3}
\begin{tabular}{@{}llll@{}}
\toprule
 & Spherically symmetric & Kerr (null framework) \\
\midrule
$\sigma_1$ & $r\sqrt{f}$ & $\sqrt{\Sigma}$ \\
$\sigma_2$ & $r$         & $\sqrt{\Sigma}$ \\
$\Box$ & $(r^{-2})\partial_r(r^2 f\,\partial_r)$
       & $\partial_r^2+(2r/\Sigma)\partial_r$ \\
$K_1=K_2=1$ & $\Leftrightarrow R_{\mu\nu}=0$ (vacuum) &
  $K{=}1$ (algebraic identity) \\
$a\to 0$ limit & --- & $\sqrt{\Sigma}\to r$, recovers left column \\
\bottomrule
\end{tabular}
\caption{Comparison of the $K{=}1$ framework in the spherically
symmetric and Kerr settings.
In Kerr, $\sigma_1=\sigma_2=\sqrt{\Sigma}$; the framework has a
single eigenvalue rather than two.}
\label{tab:comparison}
\end{table}

%───────────────────────────────────────────────────────
\section{Discussion}\label{sec:discussion}

\subsection*{Connection to $R_{\mu\nu}=0$ for Kerr}

Theorem~\ref{thm:kerr} establishes $K{=}1$ for Kerr as an
algebraic identity.
We now show that, for the principal outgoing null congruence,
$K{=}1$ is in fact \emph{equivalent} to
$R_{\mu\nu}\ell^\mu\ell^\nu=0$.

\begin{proposition}[$K{=}1\Leftrightarrow R_{\mu\nu}\ell^\mu\ell^\nu=0$]
\label{prop:equiv}
For the Kerr metric with the Kinnersley null congruence
$(\ell^r=1,\,\sigma_{\scriptscriptstyle\mathrm{NP}}=0)$ and
$\sigma=\sqrt{\Sigma}$,
\[
  K=\Sigma\,\Box_{\mathrm{null}}\ln\sigma=1
  \;\Longleftrightarrow\;
  R_{\mu\nu}\ell^\mu\ell^\nu = 0.
\]
\end{proposition}

\begin{proof}
Since $\sigma=\sqrt{\Sigma}$, one has
$D\ln\sigma=\partial_r\ln\sigma=r/\Sigma=\theta_{\mathrm{exp}}/2$,
hence $D^2\ln\sigma=D\theta/2$, so
\begin{equation}\label{eq:Kraychaudhuri}
  K = \sigma^2\bigl(D^2\ln\sigma+\theta\,D\ln\sigma\bigr)
    = \tfrac{1}{2}\sigma^2(D\theta+\theta^2).
\end{equation}
The Kinnersley congruence is geodesic ($\kappa=0$) and, by
the Goldberg--Sachs theorem~\cite{GoldbergSachs1962}
(Kerr is Petrov type~D), shear-free
($\hat\sigma_{ab}=0$).
The Raychaudhuri equation~\cite{Raychaudhuri1955} therefore reduces to
\begin{equation}\label{eq:raychaudhuri}
  D\theta = -\tfrac{1}{2}\theta^2
    + \omega_{ab}\omega^{ab}
    - R_{\mu\nu}\ell^\mu\ell^\nu,
\end{equation}
where the twist is
$\omega_{ab}\omega^{ab}=2\omega_{\scriptscriptstyle\mathrm{NP}}^2
=2\bigl(a\cos\theta/\Sigma\bigr)^2$,
with $\omega_{\scriptscriptstyle\mathrm{NP}}=\lvert\mathrm{Im}\,
\rho_{\scriptscriptstyle\mathrm{NP}}\rvert
=a\cos\theta/\Sigma$~\cite[Sec.~56]{Chandrasekhar1983}.
Substituting~\eqref{eq:raychaudhuri} into~\eqref{eq:Kraychaudhuri}:
\[
  K = \tfrac{1}{2}\sigma^2\!\left[
    \tfrac{1}{2}\theta^2+\omega_{ab}\omega^{ab}
    -R_{\mu\nu}\ell^\mu\ell^\nu\right].
\]
The \emph{kinematic identity} (verified algebraically,
using $\theta=2r/\Sigma$ and $\omega_{ab}\omega^{ab}=2a^2\cos^2\!\theta/\Sigma^2$):
\begin{equation}\label{eq:kinematic}
  \tfrac{1}{2}\theta^2+\omega_{ab}\omega^{ab}
  = \frac{2r^2}{\Sigma^2}+\frac{2a^2\cos^2\!\theta}{\Sigma^2}
  = \frac{2(r^2+a^2\cos^2\!\theta)}{\Sigma^2}
  = \frac{2}{\Sigma}=\frac{2}{\sigma^2},
\end{equation}
gives $K=1-\tfrac{1}{2}\sigma^2 R_{\mu\nu}\ell^\mu\ell^\nu$,
hence $K=1\Leftrightarrow R_{\mu\nu}\ell^\mu\ell^\nu=0$.\qed
\end{proof}

\noindent
Proposition~\ref{prop:equiv} captures the vacuum condition
in the outgoing null direction: $K{=}1$ detects whether
$R_{\mu\nu}\ell^\mu\ell^\nu$ vanishes.
By the time-reversal symmetry $t\to -t,\,\varphi\to -\varphi$
of the Kerr metric (which maps $\ell^\mu\to n^\mu$ and leaves
$R_{\mu\nu}$ unchanged), the same conclusion holds for the
ingoing congruence: $K{=}1\Leftrightarrow R_{\mu\nu}n^\mu n^\nu=0$.
In Newman--Penrose notation this gives
$\Phi_{00}=\Phi_{22}=0$.

The full equivalence $K{=}1\Leftrightarrow R_{\mu\nu}=0$
(all components) is a stronger statement.
The Kerr--de~Sitter metric provides a concrete obstruction:
it is Petrov type~D with $\Phi_{00}=\Phi_{22}=0$, yet
$\Phi_{11}\neq 0$~\cite[Sec.~57]{Chandrasekhar1983},
so $\Phi_{00}=\Phi_{22}=0$
alone does not imply $R_{\mu\nu}=0$.
Completing the equivalence requires controlling the
remaining Newman--Penrose Ricci components
$\Phi_{11},\Phi_{01},\Phi_{02},\Phi_{12}$, and remains open.

\subsection*{Why $\sqrt{\Sigma}$?}

The computation reveals why $\sqrt{\Sigma}$ is the correct
eigenvalue for Kerr.
The cancellation in the proof of Theorem~\ref{thm:kerr} is:
\[
  \underbrace{\frac{a^2\cos^2\!\theta-r^2}{\Sigma^2}}_{\partial_r^2\ln\sigma}
  +\underbrace{\frac{2r}{\Sigma}}_{\theta_{\mathrm{exp}}}
  \cdot\underbrace{\frac{r}{\Sigma}}_{\partial_r\ln\sigma}
  =\frac{1}{\Sigma}.
\]
The negative $-r^2/\Sigma^2$ from the second derivative is
exactly cancelled by the positive $2r^2/\Sigma^2$ from the
expansion term.
This cancellation is structural: it requires
$\partial_r\ln\sigma=r/\Sigma$, which uniquely (up to an
additive constant) fixes $\sigma^2\propto\Sigma$.
Any other choice of $\sigma^2$ would leave a residual
proportional to the mismatch with $\Sigma$.

\subsection*{Perturbative analysis for radial $\sigma_2$}

For completeness: if one insists on a
\emph{purely radial} $\sigma_2(r)$ for Kerr using the area-weighted
operator
\begin{equation}\label{eq:boxA}
  \Box_A\phi := \frac{1}{A}\,\partial_r\!\bigl(\Delta\,\partial_r\phi\bigr),
  \qquad A=r^2+a^2,\quad \Delta=r^2-2Mr+a^2,
\end{equation}
the $a^2$-order correction $g(r)$ defined by
$\sigma_2=r(1+a^2 g+O(a^4))$ satisfies
\begin{equation}
  \frac{d}{dr}\!\bigl[\Delta_0\,g'\bigr]+2g=\frac{2}{r^2},
  \qquad \Delta_0=r^2-2Mr.
\end{equation}
Substituting $t=(r-M)/M$ shows this is an ODE with a
regular singular point at $r=2M$ whose indicial equation
has a double root $\rho=0$; by the Frobenius method, the
homogeneous solutions near $r=2M$ form a basis where one is
a power series in $(r-2M)$ (convergent for $|r-2M|<2M$) and
the other acquires a factor $\ln(r-2M)$.
Neither is a polynomial, and the particular solution contains
$\int\!\ln(r/(r{-}2M))/r^2\,dr$, with no closed form.
This confirms that the $a^2$-order perturbative correction
has no closed-form expression.
More fundamentally, the condition $K_2{=}1$ requires
$\partial_r\ln\sigma = r/\Sigma$ (Remark~\textit{Why~$\sqrt{\Sigma}$}),
which forces $\sigma^2\propto\Sigma=r^2+a^2\cos^2\!\theta$:
no function of $r$ alone, however chosen, can reproduce the
$\theta$-dependence of $\Sigma$.

\bigskip
\noindent\textbf{Statements and Declarations}

\medskip\noindent\textbf{Funding.}
No funding was received for conducting this study.

\medskip\noindent\textbf{Competing interests.}
The author has no relevant financial or non-financial interests
to disclose.

\medskip\noindent\textbf{Data availability.}
No datasets were generated or analysed during the current study.

\begin{thebibliography}{2}

\bibitem{K1}
Y.\,Y.\,N.~Li,
$K{=}1$ Chronogeometrodynamics: Lorentzian Geometry from
Information Time.
Preprint, Zenodo (2026), \url{https://doi.org/10.5281/zenodo.19469558}

\bibitem{Raychaudhuri1955}
A.~K.~Raychaudhuri,
Relativistic Cosmology~I.
\textit{Phys.\ Rev.}\ \textbf{98}, 1123 (1955).

\bibitem{GoldbergSachs1962}
J.~N.~Goldberg and R.~K.~Sachs,
A theorem on Petrov types.
\textit{Acta Phys.\ Polon.\ Suppl.}\ \textbf{22}, 13 (1962).

\bibitem{Chandrasekhar1983}
S.~Chandrasekhar,
The Mathematical Theory of Black Holes.
Oxford University Press, Oxford (1983).
(Sec.~56: Newman--Penrose spin coefficients for the Kerr metric.)

\end{thebibliography}

\end{document}
